% !TEX root = ../../../main.tex
% Appendix entry: interrogation.event_shape -- "Event shape" (Appendix A).
\entryfield{Algorithm}
The adapter (\code{_run} in \code{Adapters/interrogation_event_shape.py}) calls
\code{measure_events} in \code{Working/interrogation/event_shape.py}, which calls
\code{measure_event} once per event. The anatomy reuses the drop detector's own rules:
\code{find_trough} (\code{Working/Detection/drop_motifs/detect.py}), \code{refine_onset} and
\code{walk_back_to_shoulder} (\code{detect5.py}) and \code{fall_gradients}
(\code{gradients.py}).

For each event with window $[s, e)$ (clipped to the span):
\begin{enumerate}
\item \emph{Orient.} Decide drop or spike (\emph{polarity}; \texttt{auto} calls it a drop when
  the steepest fall, by \code{np.gradient}, is at least as steep as the steepest rise). Form
  the oriented trace $z = o \cdot c \cdot x$ from $s$ up to the next event's start (or the
  span end), with $o = +1$ for a drop and $-1$ for a spike and $c$ the factor to millivolts
  (below). Every event is then measured as a fall.
\item \emph{Refuse} the window if it has fewer than 3 samples or any non-finite sample:
  every measure NaN.
\item \emph{Anatomy}, on the window $w = z[s{:}e]$, unless the input already carries
  \code{onset_idx} and \code{extremum_idx} columns, in which case those are the onset and
  extremum and the search is skipped:
  \begin{itemize}
  \item steepest: $k^* = \arg\min \nabla w$; if $\nabla w[k^*] \ge 0$ there is no fall;
  \item knee (\code{find_trough}): walking forward from $k^*$, the first of three
    consecutive samples whose slope is shallower than \emph{knee\_frac}$\,\cdot \nabla w[k^*]$;
  \item extremum: the lowest sample of $w$ between $k^*$ and the knee;
  \item onset: \code{refine_onset}, the same knee rule read backwards from $k^*$, then, with
    \emph{walk\_onset\_back}, \code{walk_back_to_shoulder}: step back while the previous
    sample is at least as high, never before the window start.
  \end{itemize}
\item \emph{Depth} $D = w[\text{onset}] - w[\text{ext}]$; $D \le 0$ means no fall: every
  measure NaN.
\item \emph{Measures}, with interpolated crossings between samples and $f_s$ the sample
  rate:
  \begin{description}
  \item[\code{event_amplitude_mv}] $D$: a drop's depth, a spike's height.
  \item[\code{precursor_height_mv}] $w[\text{onset}] - \min w[0..\text{onset}]$: how far the
    onset stands above the lowest point before it in the window (a drop's preceding rise, a
    spike's preceding dip).
  \item[\code{event_width_s}] $(\text{ext} - \text{onset})/f_s$.
  \item[\code{recovery_idx}, \code{recovery_time_s}] the first crossing after the extremum
    of $w[\text{ext}] + \text{\emph{recovery\_frac}}\cdot D$, searched past the window's end
    but never into the next event's window and never further than
    \emph{recovery\_max\_mult} event widths; recovery time is extremum to that crossing.
  \item[\code{duration_s}] onset to the recovery crossing.
  \item[\code{fwhm_s}] full width at half maximum: the half level is $w[\text{onset}] - D/2$;
    from its last crossing on the way down to its first on the way back (same search
    bound).
  \item[\code{rise_time_s}] spikes only (in the recorded frame): from the crossing of
    $f D$ to that of $(1-f) D$ below the onset level on the way to the extremum,
    $f$ = \emph{rise\_time\_frac} (0.1 gives the 10--90\,\% rise time). NaN for a drop,
    never 0.
  \item[\code{max_slope_mv_s}] the steepest value of $\nabla w \cdot f_s$ over
    $[\text{onset}, \text{ext}]$.
  \item[\code{onset_slope_mv_s}] $\nabla w \cdot f_s$ at the onset sample.
  \item[\code{chord_slope_mv_s}] $(w[\text{ext}] - w[\text{onset}]) / \text{width}$.
  \item[\code{peakedness}] max slope / chord slope: 1 for a straight fall, larger for a
    front-loaded one.
  \item[\code{span_ptp_mv}] peak to peak over the whole window.
  \end{description}
\item \emph{Back to the recorded frame.} Signed quantities (the three slopes) are multiplied
  by $o$, and by $-1$ again when \emph{upstream\_inverted}; \code{polarity} is $-1$ for a
  drop and $+1$ for a spike in the recorded frame. Amplitudes and widths stay positive.
\end{enumerate}

\emph{The rose.} \code{event_rose} hands the oriented snippets to \code{gradients.rose_data}
(field \code{max_slope_mv_s}), split by an upstream \code{group} column when there is one.
Every event is drawn as a fall, so a spike's rise is drawn as its mirror. What 45\textdegree{}
means is set by \emph{rose\_scale}; on \texttt{raw} it is \emph{rose\_reference\_mv\_s}, or,
at 0, the Library's stored rose reference (\code{Working/library/rose_reference.py}), or, when
none is stored, these events' own pooled median, said in the rose's text.

\entryfield{Assumptions}
\begin{itemize}
\item \emph{Units.} The block is unit-aware: \code{_run} takes the recording's declared
  \code{units} and \code{Working/units.py}'s \code{read_as_mv} gives the factor $c$ (V
  $\times$1000, mV $\times$1, µV $\times$0.001). An undeclared unit is read as volts and
  \code{meta.units_note} says so. Amplitudes are in mV, slopes in mV/s (the sample rate is
  applied once), widths and times in seconds.
\item \emph{Sampling rate.} Slopes are $\nabla x \cdot f_s$, so they are comparable across
  sampling rates; the knee rule counts samples (three in a row), so at higher rates it is a
  shorter time.
\item \emph{Baseline.} The detrended trace is the right input: the slow drift of a raw trace
  adds to every slope (\code{fall_gradients}' docstring). The block measures whatever signal
  the chain hands it.
\item \emph{Missing values.} Any NaN inside an event's window: no measures for that event
  (counted, see below). NaNs beyond the window stop the recovery and FWHM searches, since a
  comparison with NaN is false: recovery is then NaN.
\item \emph{Edges.} A window running past the span is clipped to it. The recovery search of
  the last event may run to the span's end.
\item \emph{Indices} (\code{onset_idx}, \code{extremum_idx}, \code{recovery_idx}) are in
  the span-relative frame, like the span starts; \code{recovery_idx} can be fractional.
\end{itemize}

\entryfield{Outputs}
A \texttt{SpanSet} with the spans, labels and scores unchanged and a feature table with one
row per event: \code{polarity}, \code{onset_idx}, \code{extremum_idx}, \code{recovery_idx},
\code{event_amplitude_mv}, \code{precursor_height_mv}, \code{event_width_s},
\code{duration_s}, \code{fwhm_s}, \code{recovery_time_s}, \code{rise_time_s},
\code{max_slope_mv_s}, \code{onset_slope_mv_s}, \code{chord_slope_mv_s}, \code{peakedness},
\code{span_ptp_mv}, followed by any upstream columns (one of the same name is replaced).

\code{meta}: \code{rules} (the printed rule for every measure, including the unit), the
\code{rose} payload, \code{n_events}, \code{n_no_fall}, \code{n_not_recovered},
\code{n_no_fwhm}, \code{n_no_rise}, \code{measures}, \code{units}, \code{to_mv} and
\code{units_note}.

\entryfield{Cost}
Each event's work reaches at most to the next event's start, so the total is $O(n)$ in the
span length plus the knee walks. The per-event loop is plain Python. The rose is a histogram
over the events.

\entryfield{Failure modes}
\begin{itemize}
\item \emph{Counted under the wrong name.} A window holding one NaN, a window of fewer than
  3 samples and a window with no fall are all counted in \code{n_no_fall} (and
  \code{n_no_rise}); the reason (\code{why} in \code{measure_event}) is not passed to the
  output.
\item \emph{\texttt{auto} polarity} called 50 of 410 seed drops spikes, because their windows
  carry minutes of context with a rise as fast as the fall (module docstring).
\item \emph{The wrong fall.} In a wide window the steepest sample can belong to a
  neighbouring event: on 360 stored snippets the search matched the detector's trough within
  3 samples on 345 and locked onto a steeper neighbouring fall on 10 (module docstring).
  Upstream \code{onset_idx}/\code{extremum_idx} columns avoid this.
\item \emph{Square bottoms.} On a flat floor the extremum is the floor's first sample, and
  the onset can land well inside a slow dip: on a synthetic square dip the onset landed
  0--230\,s after the dip's start (\code{Adapters/interrogation_event_families.py},
  docstring).
\item \emph{No recovery.} An event that never climbs back within the search bound has NaN
  recovery, duration and often FWHM; counted.
\item \emph{A unit not declared} is read as volts: a recording stored in mV then reads
  1000 times too large, with the note printed.
\end{itemize}

\entryfield{Validation}
\code{tests/test_adapter_event_shape.py} builds a drop by hand (10\,Hz, volts: a rise to
+2\,mV, a fall to $-8$\,mV over 2\,s, a floor, a linear recovery) and checks every measure
against values worked out on paper: onset 50, extremum 70, amplitude 10\,mV, precursor
2\,mV, width 2\,s, FWHM 4.5\,s, recovery 3.5\,s, duration 5.5\,s, max and chord slope
$-5$\,mV/s, peakedness 1, and the exact onset slope. It also checks a spike as the mirror of
a drop, the 10--90\,\% rise time of a spike, NaN rather than 0 for a drop's rise time, the
recovery search stopping at the next event, upstream anchors, and upstream inversion.
\code{tests/test_units_reach_the_blocks.py} checks true millivolts on a channel stored in mV.
The module docstring reports agreement with the detector's own depths on the 410 seed events
when the store's anchors are used. No test compares the measures with hand-marked events on a
fungal recording.

\entryfield{References}
The anatomy rules are the project's drop detector's (\code{detect.py}, \code{detect5.py});
FWHM, recovery, duration and precursor are defined in
\code{Working/interrogation/event_shape.py}. The 10--90\,\% rise time is the usual
electrophysiological convention. No external paper defines this block.

\entryfield{Tests}
\code{tests/test_adapter_event_shape.py} (24 tests, above); \code{tests/test_rose_reference.py}
(the rose's 45\textdegree{}); \code{tests/test_units_reach_the_blocks.py} (units);
\code{tests/test_webui_interrogation_shape.py} (the page's use of it); and the template
chains in \code{tests/test_template_event_features.py}, \code{test_template_event_trains.py},
\code{test_template_event_families.py}, \code{test_template_chain_cycles.py} and others run it
inside chains.
